\rhead{}
\lhead{}
\renewcommand{\headrulewidth}{0pt}
\addcontentsline{toc}{chapter}{\textbf{Resumen}}
\begin{center}
 
    \textbf{Resumen}\\
\end{center}

\noindent
MatEnergy-ML es una plataforma computacional que aplica aprendizaje automático
clásico sobre datos de Teoría del Funcional de la Densidad (DFT) para acelerar
el cribado de materiales candidatos para cátodos y ánodos de baterías de
ion-litio. En lugar de calcular propiedades cuánticas costosas para cada
composición candidata, el sistema entrena modelos de regresión (\textit{Ridge
Regression}, \textit{Random Forest}, \textit{Gradient Boosting}) y clasificación
(\textit{Random Forest}) sobre un conjunto de 151 materiales reales extraídos de
Materials Project, prediciendo energía sobre el casco convexo, energía de
formación, \textit{band gap} y estabilidad termodinámica binaria a partir de
descriptores composicionales calculados con \texttt{pymatgen}. La plataforma
implementa una arquitectura limpia por capas (FastAPI, PostgreSQL, React), con
autenticación JWT, control de acceso basado en roles, verificación de
integridad de artefactos mediante SHA-256, y un sistema de \textit{ranking}
transparente y no generativo para priorizar candidatos según su relevancia
para aplicaciones de almacenamiento de energía. Los resultados muestran que
\textit{Gradient Boosting} alcanza un MAE de prueba de 0.0186~eV/átomo para
energía sobre el casco convexo, mientras que \textit{Random Forest} predice la
energía de formación con $R^2=0.869$; el clasificador de estabilidad alcanza
F1-macro$=0.583$ sobre el conjunto de prueba. Se discuten las limitaciones del
tamaño del conjunto de datos y del alcance puramente composicional de los
descriptores, así como la diferencia entre estabilidad termodinámica y
desempeño electroquímico real de una batería.

\vspace{0.6cm}
\noindent
\textbf{Palabras clave:} aprendizaje automático, ciencia de materiales
computacional, baterías de ion-litio, DFT, cribado de materiales.
